\begin{table}[htbp]\centering\footnotesize
\caption{Razones por persona: cada una con su denominador declarado (pesos de cada año)}\label{cua:percap}
\begin{tabular}{lrrr}
\toprule
Razón & 2026 & 2027 & Var. real (\%) \\
\midrule
Perímetro pensionario / población 65 y más & 139.670 & 144.665 & $-$0,4 \\
FONE / población en edad escolar básica & 21.494 & 24.102 & +7,8 \\
Ramos 12 y 56 / población total & 1.781 & 1.967 & +6,2 \\
\midrule \multicolumn{4}{l}{\emph{El mismo numerador con dos denominadores distintos, para que se vea la diferencia}} \\
FONE 2027 / población en edad escolar básica 2027 & --- & 24.102 & --- \\
FONE 2027 / matrícula de básica, ciclo 2025--2026 & --- & 26.564 & --- \\
\bottomrule
\end{tabular}
\\[2pt]\begin{minipage}{0.92\textwidth}\scriptsize \textbf{Cada renglón declara denominador, año y cobertura.} Población: CONAPO, a mitad de año, nacional. Matrícula: SEP, ciclo 2025--2026, modalidad escolarizada, nacional. \textbf{Matrícula no es población en edad escolar}: los dos denominadores difieren en 10,2~\% y mueven la cifra por persona en esa misma proporción. \textbf{El primer renglón es un indicador de presión demográfica, NO una pensión media}: el denominador incluye a quien no recibe y el numerador puede incluir pensiones de menores de 65 años.\end{minipage}
\end{table}
